\noindent\textbf{Elección del procedimiento.} Usar el convenio de función de corriente $u=\partial_y\psi$, $v=-\partial_x\psi$. Para el polinomio dado se obtiene $u=-6Kxy$, $v=-3K(x^2-y^2)$. Primero se derivan las funciones, después se evalúa el punto.

\begin{enumerate}

\item La derivada de $-3xy^2$ respecto a $y$ aporta $-6xy$.

\[\fitmath{u=-6\cdot K\cdot x\cdot y}\]

\[\fitmath{u=-6\cdot (0{,}2)\cdot (1)\cdot (1)\approx -1{,}2\,\mathrm{m/s}}\]

\item El signo menos pertenece a la definición de $v$.

\[\fitmath{v=-3\cdot K\cdot \left({x}^{2}-{y}^{2}\right)}\]

\[\fitmath{v=-3\cdot (0{,}2)\cdot \left({(1)}^{2}-{(1)}^{2}\right)\approx 0\,\mathrm{m/s}}\]

\item La celeridad es el módulo del vector velocidad.

\[\fitmath{\mathrm{speed}=\sqrt{{u}^{2}+{v}^{2}}}\]

\[\fitmath{\mathrm{speed}=\sqrt{{(-1{,}2)}^{2}+{(0)}^{2}}\approx 1{,}2\,\mathrm{m/s}}\]

\end{enumerate}

\noindent\textbf{Contraste y comparación.} $u_x=-6Ky$ y $v_y=6Ky$ cancelan la divergencia; $v_x-u_y=-6Kx+6Kx=0$. El flujo prescrito es incompresible e irrotacional.
